\documentclass[11pt]{article}
\usepackage[margin=1in]{geometry}
\usepackage{amsmath,amssymb,booktabs,hyperref}
\usepackage[numbers]{natbib}
\newcommand{\todo}[1]{\textcolor{red}{[TODO: #1]}}
\usepackage{xcolor}
% RULE: every number in this paper is \input from generated/ (scripts/make_tables.py).
\title{Open Problems in World Models and Interpretability}
\author{\todo{authors}}
\begin{document}
\maketitle

\begin{abstract}\todo{}\end{abstract}

% Template per problem: Statement / Why it matters / Current evidence (cite) / What would count as progress (a test).
\section{Defining the object of study}
\paragraph{P1: What makes a representation a world model?} Decodability is necessary, not sufficient. \todo{}
\paragraph{P2: Ground truth beyond toy worlds.} \todo{}

\section{Representations}
\paragraph{P3: Memory vs.\ inference of latent state in partially observed settings.} \todo{}
\paragraph{P4: Identifiability of latent state variables.} \todo{}

\section{Dynamics}
\paragraph{P5: Locating transition computation; error accumulation in imagined rollouts.} \todo{}

\section{Use in planning}
\paragraph{P6: Does the planner's cost reflect reachability?} Evidence from the companion paper. \todo{}
\paragraph{P7: Hindsight structure in latent spaces.} \todo{}

\section{Methodology}
\paragraph{P8: Probe validity and controls as a standard.} \citep{hewitt2019control} \todo{}
\paragraph{P9: Causal evaluation at scale (video WMs, JEPAs, world action models).} \todo{}

\section{Relation to existing roadmaps}
\citep{openproblems_mi_2501_16496,wm_roadmap_2607_06401,state_of_mi_2611_08438}
\todo{Each problem above must state what is new relative to these.}

\bibliographystyle{plainnat}
\bibliography{../refs}
\end{document}
